\subsection{例子}

在本条中，我们假设 X 是纯有限 n 维的。我们还假设它装备了一个 \(C^{r+1}\) 类流形结构，它与给定的 \(C^{r}\) 类结构相容。

14.4.1（“无穷小变换”）。设 \(\tau\) 是一个关于同构的有限维向量函子，为 \(C^r\) 类；我们假设对每个有限维向量空间 V，\(\tau(V)\) 是有限维的。用 \(E_t\) 表示向量丛 \(\tau(T(X))\)；它是 \(C^r\) 类的。

设 \(\xi\) 是 X 上的一个 \(C^r\) 类向量场。存在一个 \(E_{\tau} \to E_{\tau}\) 型、阶 \(\leqslant1\) 的微分算子 D，使得

\[
\mathbf {D} _ {\mathrm{U}} (s) = \theta_ {\xi}. s
\]

对 X 的每个开集 U 及每个 \(s \in \mathcal{S}_{\mathbb{E}}^{r}(\mathrm{U})\) 成立（参见 8.4.3）。这样的算子 D 是唯一的；我们把它记作 \((\theta_{\xi})_{\tau}\) 或简称 \(\theta_{\xi}\)。它的符号是向量丛的截面 \(\xi \otimes \mathrm{Id}_{\mathbb{E}_{\tau}}\)，该向量丛即

\[
\mathrm{S} ^ {1} (\mathrm{E} _ {\tau}, \mathrm{E} _ {\tau}) = \mathrm{T} (\mathrm{X}) \otimes \mathscr {L} (\mathrm{E} _ {\tau}; \mathrm{E} _ {\tau}).
\]

若 V 是一个有限维 K-向量空间，令

\[
\tilde {\tau} (\mathbf {V}) = \mathscr {L} (\tau (\mathbf {V}); \mathrm{Alt} ^ {n} (\mathbf {V}; \mathbf {K})) = \tau (\mathbf {V}) ^ {*} \otimes \bigwedge^ {n} \mathbf {V} ^ {*},
\]

且若 \(u\colon V_{1}\to V_{2}\) 是一个同构，则通过转移结构定义 \(\tilde{\tau}(u)\colon\tilde{\tau}(V_{1})\to\tilde{\tau}(V_{2})\)。这样就得到一个有限维的向量函子 \(\tilde{\tau}\)。丛 \(E_{\tilde{\tau}}=\tilde{\tau}(T(X))\) 以显然的方式与丛 \((E_{\tau})^{\sim}\) 恒同；\((\theta_{\xi})_{\tau}\) 的转置是 \(-( \theta_{\xi})_{\tilde{\tau}}\)，且 14.3.4 中的公式 (1) 取如下形式

\[
\langle \theta_ {\xi}. u, v \rangle + \langle u, \theta_ {\xi}. v \rangle = d (i (\xi) (u. v)).
\]

14.4.2（“外微分”）。对每个整数 \(p \geqslant 0\)，令

\[
\Omega^ {p} = \operatorname{Alt} ^ {p} (\mathrm{T} (\mathrm{X}); \mathrm{K} _ {\mathrm{X}}).
\]

存在一个 \(\Omega^{p} \to \Omega^{p+1}\) 型、阶 \(\leqslant 1\) 的微分算子 D，使得

\[
\mathrm{D} _ {\mathrm{U}} (\omega) = d \omega
\]

对 X 的每个开集 U 及每个 \(\omega\in\mathcal{S}_{\Omega^{p}}^{r}(U)\) 成立。这样的算子 D 是唯一的；把它记作 \(d\)（或记作 \(d_{p}\)，如果想指明整数 \(p\)）。

它的符号是 \(S^1(\Omega^p, \Omega^{p+1})\) 的元素，其求法如下：首先有

\[
\begin{array}{r l} \mathrm{S} ^ {1} (\Omega^ {p}, \Omega^ {p + 1}) & = \mathrm{T} (\mathrm{X}) \otimes \mathcal {L} (\Omega^ {p}; \Omega^ {p + 1}) \\ & = \mathcal {L} (\Omega^ {1}; \mathcal {L} (\Omega^ {p}; \Omega^ {p + 1})) = \mathcal {L} (\Omega^ {1}, \Omega^ {p}; \Omega^ {p + 1}), \end{array}
\]

其中 \(\mathcal{L}(\Omega^{1},\Omega^{p};\Omega^{p+1})\) 表示从 \(\Omega^{1}\times\Omega^{p}\) 到 \(\Omega^{p+1}\) 的双线性映射组成的丛。而外积 \((\alpha,\beta)\mapsto\alpha\wedge\beta\) 定义了后一丛的一个典范截面；这个截面就是 \(d_{p}\) 的符号。

为确定 \(d_{p}\) 的转置，首先注意外积定义了一个配对 \(\Omega^{p} \times \Omega^{n-p} \to \Omega^{n} = \Omega\)，借助它把 \((\Omega^{p})^{\sim}\) 与 \(\Omega^{n-p}\) 恒同。以同样方式把 \((\Omega^{p+1})^{\sim}\) 与 \(\Omega^{n-p-1}\) 恒同。在这些约定下，\(d_{p}\) 的转置是 \((-1)^{p+1}d_{n-p-1}\)，而相应的 \((\Omega^{p}, \Omega^{n-p-1}) \to \Omega_{1} = \Omega^{n-1}\) 型 Green 算子就是外积（在每个纤维上）。

14.4.3（“Laplace 算子”）。取 K = R，r = ∞，X = R \(^{n}\)，E = F = K \(_{X}\)；用 \(\xi^{i}\) (1 ≤ i ≤ n) 表示 X 上的坐标函数；令 D \(_{i}\) = ∂/∂ξ \(^{i}\)，以及 L = ∑ \(_{i=1}^{n}\) D \(_{i}^{2}\)。算子 L 是一个阶 ≤ 2 的标量微分算子。若把 Ω 与 K \(_{X}\) 恒同（借助标架 ω = dξ \(^{1}\) ∧ ⋯ ∧ dξ \(^{n}\)），则有 \(\tilde{E}\) = \(\tilde{F}\) = K \(_{X}\)（参见 14.3.3 的例子）且 \(^{t}\) L = L。

对 \(1 \leqslant i \leqslant n\)，令 \(\omega_{i} = (-1)^{i-1} d\xi^{1} \wedge \cdots \wedge d\xi^{i-1} \wedge d\xi^{i+1} \wedge \cdots \wedge d\xi^{n}\)。诸 \(\omega_{i}\) 构成 \(\Omega_{1} = \Omega^{n-1}\) 的一个标架。若 \(f \in \mathcal{C}^{1}(X)\)，用 \(\text{grad}(f)\) 表示截面 \(\sum_{i=1}^{n} D_{i}(f)\omega_{i}\)（\(\Omega_{1}\) 的截面）。存在 L 的一个 Green 算子 G，使得有

\(\mathbf{G}(u,v) = v.\operatorname{grad}(u) - u.\operatorname{grad}(v)\) 对 X 的每个开集 U 及 \(u,v \in \mathcal{C}^{1}(\mathbf{U})\) 成立。

特别地，设 A 是 \(\mathbf{R}^{n}\) 的一个紧块。给 \(\mathbf{R}^{n}\) 定向，给 A 装备由 \(\mathbf{R}^{n}\) 的定向诱导的定向，并给 \(\partial\mathbf{A}\) 装备相应的定向（11.2.1）。我们有

\[
\int_ {\mathbf {A}} (\mathbf {L} u). v \omega - \int_ {\mathbf {A}} u. (\mathbf {L} v) \omega = \int_ {\partial \mathbf {A}} (v. \operatorname{grad} (u) - u. \operatorname{grad} (v))
\]

其中 u、v 取遍 \(\mathcal{C}^{1}(X)\)。

